\section{Introduction}
\label{sec:introduction}
An agentic platform turns one user request into many executions. An orchestrator delegates to specialists; specialists fan out to workers on threads; a run pauses for approval and resumes in another process; a cache answers a repeated prompt; a retriever injects documents into context; a tool is exported to a different framework. The request, however, arrived with controls attached: who the caller is, which tenant owns the data, which scopes were granted, how much may be spent, how many calls per minute are allowed, and what to do if the policy service is down. Each boundary above is a place where those controls can silently disappear.

This is not hypothetical. In a source review of \method{} before this work, request deadlines, cancellation and deny-all tool policies were ignored at the public API in both of its runtimes; a viewer's request routed through a native supervisor executed an admin-only tool while the LangGraph runtime correctly denied it; a human approval overrode an external policy engine that was never consulted; and a prompt cache ignored tool schemas (\S\ref{sec:conformance}). Public evidence shows the same failure class across the ecosystem: on AgentGovBench~\cite{agentgovbench}, each of seven widely used frameworks and coding agents scores \AGBNativeScore{}/\AGBTotal{} when run without an external governance layer: as run, none propagates the end user's identity into an audit record, narrows delegated scope, or aggregates rate limits across subagents.

We argue that these are not independent bugs but consequences of an architectural choice. Many frameworks treat controls as \emph{arguments}: values that each layer must receive and forward. Every new boundary is then a new opportunity to forget one. \method{} instead treats controls as a property of the \emph{execution scope} of a request. Scopes compose by a meet operation, so nesting can only narrow them; live handles (budgets, cancellation tokens) are kept out of persisted state; and a single deny-first admission point decides and audits every tool call. We call the resulting specification the \emph{Governed Execution Contract} (\gec{}).

\paragraph{Research questions.}
\textbf{RQ1} (foundations): can the control semantics an enterprise platform needs be stated as a small algebra whose laws are provable and machine-checkable? \textbf{RQ2} (governance): does a boundary-preserving architecture meet a public governance benchmark, and which mechanisms are responsible for which outcomes? \textbf{RQ3} (measurement): do existing governance benchmarks distinguish correct enforcement from trivial behavior? \textbf{RQ4} (cost and breadth): what do the controls cost, and how does the platform perform on the other layers---orchestration, context, security filters, evaluation, task execution and scale?

\paragraph{Contributions.}
\begin{enumerate}
\item \textbf{A formal contract.} The \gec{} (\S\ref{sec:formal}) models control scopes as a meet-semilattice and states seven invariants. We prove attenuation along arbitrary delegation chains, deny-monotonicity of policy resolution, exactness of per-principal rate limits under arbitrary interleavings, decision totality under policy-source failure (including an impossibility argument for reading the fail mode from the failed source), and coordination-free sharding by principal. Property-based tests check the algebraic laws on randomly generated policies and chains.
\item \textbf{A platform architecture.} \method{} (\S\ref{sec:architecture}) realizes the contract across seven layers: declarative entry points, a think--act--critique orchestration loop, a runtime seam with a native engine and a LangGraph engine, a governed tool gateway, a multi-tenant control plane, an LLM gateway, and a budgeted context layer, with evaluation and observability as cross-cutting rails. We describe the design decisions that differ from common practice and why.
\item \textbf{State-of-the-art governance results with mechanism attribution.} \method{} passes \AGBFoundry{}/\AGBTotal{} AgentGovBench scenarios in \AGBRepeatsPerfect{}/\AGBRepeats{} repetitions, above every published configuration (best \AGBBestPublished{}/\AGBTotal{}) and above the two closest open-source systems run under our harness (Agent Governance Toolkit at most \AGTBestOfficial{}, Bounded Agents \APCOfficial{}); on 24 scenarios written blind to our code it scores \IndepFoundry{}/24 against \AGTBestIndep{} and \APCIndep{}. Single-control ablations (\AblationCount{}) attribute \AttributedOfficial{} scenarios to specific mechanisms (\S\ref{sec:governance}).
\item \textbf{A measurement finding and remedy.} Extending trivial-agent sanity checks~\cite{zhu2025abc} to an enforcement benchmark: Constant policies pass most individual scenarios (\VacuousAny{}/\AGBTotal{} are passed by at least one trivial runner); a deny-everything runner scores \AGBDenyAll{}/\AGBTotal{}; and \OfficialBlindControls{} controls (\OfficialBlindList{}) can be removed without losing a point. Published runs corroborate this: \ACPDeclaredNarrowingPassed{} of \ACPDeclaredNarrowing{} governed runs pass the task-narrowing scenario while declaring that they do not implement task narrowing. We contribute eight discriminating scenarios and three anti-vacuity checks (\S\ref{sec:vacuity}).
\item \textbf{A 360-degree evaluation} across orchestration overhead, governance overhead, thread and process scaling, long-context assembly, public injection and hallucination datasets, and matched WorkBench task execution (\S\ref{sec:breadth}), with negative results kept: the default marker-based injection filter detects \InjRegexRecall{} of held-out attacks (we add a trained gate that detects \InjNBRecall{}), keyword retrieval fails paraphrased queries, and task accuracy does not improve significantly over the benchmark's reference loop.
\end{enumerate}

\paragraph{Scope of claims.}
We claim state of the art only where we measure it against published results under the same scorer: governance on AgentGovBench, and coordination overhead against LangGraph on identical topologies. We do not claim task-level superiority on WorkBench, where the best published model run reaches \WBFrontierCorrect{}/690 with a stronger model and a different harness, nor multi-agent \emph{quality} advantages, which require model-driven benchmarks we did not run. The control plane was developed during this study with knowledge of AgentGovBench's categories; we treat this as a threat to validity and address it with ablations, self-authored discriminating scenarios, and anti-vacuity checks rather than by assertion (\S\ref{sec:threats}).

\paragraph{Artifact.}
Code, evidence and this manuscript are at \url{https://github.com/arghya05/agent-foundry}. Every number in this paper is generated by \code{research-paper/build\_evidence.py} from raw files under \evpath{review/evidence/}; the main campaign was run from commit \code{\CampaignCommit{}} and produced \CampaignFiles{} files with SHA-256 digests in its manifest.
